\subsection{Adaptive Investigation and the Validity of Scientific Evidence}
\label{sec:adaptive-evidence}

An agent that chooses experiments changes the evidence available for subsequent inference. Its scientific contribution therefore depends on the relation between the decision objective, the distribution of acquired observations, and the claim made from those observations. The distinction between property optimization and hypothesis discrimination in Section~\ref{sec:workflow-delegation} has an important consequence: a policy can become better at finding a desirable object while learning less about other parts of the system. Exploration within Bayesian optimization can support the search for an optimum without being designed to identify a mechanism or characterize the entire response surface.

This divergence is observable in a computational materials study. Yoshida and colleagues evaluated linked autonomous searches for magnetization, Curie temperature, and spin polarization using a fixed dataset of 16,908 alloys calculated by density functional theory. Transferring learned representations improved some search outcomes, yet the Curie-temperature search preferentially acquired high-temperature examples, and predictive performance on the remaining materials, which included a substantial fraction with low or zero Curie temperatures, temporarily deteriorated~\cite{asd0513}. Better candidate acquisition and better prediction throughout the search space were therefore separate achievements. This result concerns a computed candidate library, rather than networked physical laboratories. It also shows why the scientific value of a policy must be assessed against its intended question: locating a high-performing material and estimating a general composition--property relation require different coverage.

Mechanism identification imposes an additional constraint. If competing hypotheses predict the same observation distribution for every accessible action and history, a policy restricted to those actions cannot distinguish them. More iterations can improve precision without resolving this observational equivalence. The available interventions, measured quantities, and hypothesis representation jointly determine what is identifiable. In the Robot Scientist, growth-rescue assays were informative because different gene--enzyme assignments implied different outcomes within the supplied pathway model~\cite{king2004}. The inference remained conditional on that representation. In the AI Theorist study, proposed polarization-dependent spectra identify an additional measurement relevant to testing the excitonic assignments; the corresponding experiment had not been performed~\cite{asd1081}. Improving a fit to existing observations and acquiring an observation that discriminates between explanations thus address different scientific needs. For an agent, recognizing when another measurement or model class is required is a distinct decision from continuing the current search.

\begin{samepage}
Adaptive acquisition does not, by itself, invalidate inference. Let $D_t=\{(a_s,y_s)\}_{s=1}^{t}$ record actions and outcomes, $\pi$ denote the action-selection policy, and $\theta$ index the observation model. For a fixed number $T$ of actions, when the policy depends on $\theta$ only through the recorded history, the joint distribution factorizes as
\begin{equation}
p_\theta(D_T)=\prod_{t=1}^{T}\pi(a_t\mid D_{t-1})\,p_\theta(y_t\mid a_t,D_{t-1}).
\label{eq:adaptive-evidence}
\end{equation}
\end{samepage}
The policy factors contain no additional dependence on $\theta$, so a correctly specified likelihood built from the sequential observation factors can accommodate the adaptive design. This argument assumes a fully recorded decision process and a fixed horizon; it does not establish validity for arbitrary stopping rules, unrecorded selection, or model misspecification. Nor does it guarantee that conventional confidence intervals retain their nominal coverage. Hadad and colleagues show that adaptive allocation can produce biased or heavy-tailed conventional estimators and develop inference under explicit conditions on assignment probabilities and weighting~\cite{hadad2021}. Their results concern randomized adaptive experiments; they are not an automatic correction for deterministic laboratory optimizers. Where some actions are never selected, reweighting alone cannot reconstruct evidence about them without further assumptions or new observations.

Interpreting an acquired dataset consequently requires attention to how it was generated. In the sputter-epitaxy study, a random forest fitted to the initial optimization data achieved an in-sample $R^2$ of 0.92 but approximately 0.43 under five-fold cross-validation. The authors explicitly attributed the more limited predictive performance to sample size and the acquisition-biased distribution, and used the model as a local interpretable surrogate~\cite{asd0999}. A random split of such a dataset still tests prediction within the acquired distribution; it does not establish accuracy throughout the original design space. Feature importance and partial dependence summarize the fitted response under those conditions, without independently identifying a microscopic causal mechanism. The subsequent experiments using the inferred growth rule, discussed in Section~\ref{sec:domains}, supply a different kind of support: the rule guided new decisions beyond fitting the original acquisition data.

The preceding inference also depends on an adequate observation model. A residual between prediction and measurement does not, by itself, identify whether the predictive model, measurement procedure, or physical explanation requires revision. Section~\ref{sec:model-observation} examines how selected experimental systems address these different possibilities.

The endpoint determines how adaptive selection affects the final claim. For AlphaEvolve's mathematical construction, correctness rests on explicit geometric conditions and a valid implication from the finite witness~\cite{asd0063}. Adaptive search over many candidates does not weaken that implication once its premises are checked. An empirical optimum or association instead depends on measurement variability, the sampled conditions, and the relevant inferential design. Likewise, a low predicted improvement can justify ending a resource-limited optimization campaign without establishing a unique mechanism or excluding hypotheses outside the represented family. Scientific autonomy therefore includes deciding whether the available evidence answers the research question, whether a different experiment is needed, and whether the current formulation can answer it at all. This analysis extends workflow participation to the control of what can be learned from the investigation.
